%%%%%%%%%%  scan idx 46  |  folio 33  %%%%%%%%%%
\emph{ce qui équivaut à~:}
\begin{equation}\tag{1.30 bis}
\ch(\lambda_{-1}(\check N))=N^q\mathcal C(N)^{-1}.
\end{equation}

\subsection*{\S 4. Les opérations $\lambda^p(N,x)$}

Soit $A$ l'anneau des séries formelles, à coefficients dans $\Z$, en deux suites infinies d'indéterminées, notées $\lambda^iN$ et $\lambda^i x$ ($i\geq1$). On pose $\lambda^1N=N$ et $\lambda^1x=x$, et l'on introduit sur $A$ l'unique structure de $\lambda$-anneau spécial pour lequel on ait $\lambda^i(N)=\lambda^iN$ et $\lambda^i(x)=\lambda^ix$ et pour lequel les opérations $\lambda^i$ soient «~continues~» (si l'on se bornait aux \emph{polynômes} en des indéterminées $\lambda^iN$ et $\lambda^ix$, on aurait le «~$\lambda$-anneau libre engendré par $N$ et $x$~»~; notre anneau $A$ en est un complèté). On peut alors former
\begin{equation}\tag{1.31}
\lambda_{-1}(N)=\sum_{i\geq0}(-1)^i\lambda^iN,
\end{equation}
qui est une série formelle de terme constant $1$, donc inversible~; on peut donc définir sans ambiguïté un élément $\lambda^p(N,x)$ de $A$ en posant
\begin{equation}\tag{1.32}
\lambda^p(N,x)\lambda_{-1}(N)=\lambda^p(x\lambda_{-1}(N))
\end{equation}

\textbf{Théorème 1.4.} \emph{Soit $J_q$ l'idéal fermé de $A$ engendré par les $\lambda^iN$ avec $i>q$~; identifions $A/J_q$ à l'anneau des séries formelles en les $\lambda^i x$ et les $\lambda^jN$ ($1\leq j\leq q$). Alors $\lambda^p(N,x)$ réduit modulo $J_q$ est un} polynôme, \emph{isobare de poids $p$ par rapport aux variables $\lambda^i x$.}

\emph{Soient $N$ et $F$ deux espaces vectoriels de dimensions finies $q,q'$ sur un corps $k$ algébriquement clos de caractéristique $0$~; soit $N^{(p)}$ le noyau de l'application} $(a_1,\ldots,a_p)\longrightarrow a_1+\cdots+a_p$ de $N^p$ \emph{dans} $N$.
On 
